\subsection{Общая архитектура экспериментальной установки}

Экспериментальная установка реализует принцип работы SFCW-радара, описанный в
предыдущей главе, на основе радиофотонного генератора СВЧ-сигнала. Общая
структурная схема установки представлена на рис.~\ref{fig:setup-block-diagram}.

\begin{figure}[H]
    \centering
    \fbox{\parbox[c][0.25\textheight][c]{0.85\textwidth}{\centering
    Место для структурной схемы экспериментальной установки:\\[6pt]
    Двухлазерный источник $\to$ Фотодиод (гетеродинирование)
    $\to$ СВЧ-тракт 3{,}3--6{,}3\,ГГц $\to$ Передающая антенна
    $\to$ Объект $\to$ Приёмная антенна $\to$ Синхронный детектор
    $\to$ АЦП $\to$ ПК (обработка и визуализация).}}
    \caption{Структурная схема радиофотонной экспериментальной установки}
    \label{fig:setup-block-diagram}
\end{figure}

Установка состоит из пяти функциональных подсистем, взаимодействующих через
аналоговые, цифровые и оптические каналы. \textbf{Двухлазерный оптический
источник} генерирует два одномодовых оптических излучения вблизи длины волны
1550\,нм с независимой стабилизацией температуры и тока каждого лазера. Оба
луча совмещаются на высокоскоростном \textbf{фотодиоде} (оптическом смесителе),
в фототоке которого возникает компонента на разностной частоте
$f_{\mathrm{RF}} = |\nu_1 - \nu_2|$; перестройка температуры или тока одного
из лазеров при фиксированных параметрах второго реализует частотный свип в
СВЧ-диапазоне. Сформированный сигнал излучается через \textbf{антенный блок} ---
передающую и приёмную антенны, разнесённые на фиксированное расстояние в
бистатической конфигурации. Отражённый сигнал поступает в \textbf{подсистему
регистрации}: синхронный детектор выделяет синфазную ($I_k$) и квадратурную
($Q_k$) составляющие относительно опорного гетеродина, а АЦП оцифровывает
результат с разрядностью 16\,бит. Наконец, \textbf{управляющий ПК} задаёт
параметры свипа, принимает данные и выполняет обработку --- формирование
комплексного S-параметра $S_{21}(f_k)$, обратное преобразование Фурье,
построение A- и B-сканов.

\paragraph{Принцип генерации СВЧ-сигнала.}
В отличие от традиционных схем, использующих электронный генератор
качающейся частоты (ГКЧ) или синтезатор, данная установка формирует
радиочастотный сигнал оптическим способом. Два полупроводниковых лазера с
длинами волн $\lambda_1$ и $\lambda_2$ совмещаются на фотодиоде, в фототоке
которого возникает биение на частоте
\begin{equation}
    f_{\mathrm{RF}} = c\left|\frac{1}{\lambda_1} - \frac{1}{\lambda_2}\right|.
    \label{eq:beat-freq}
\end{equation}
Для лазеров C-диапазона ($\lambda \approx 1550$\,нм) сдвиг длины волны на
$\sim$0{,}025\,нм соответствует изменению разностной частоты на $\sim$3\,ГГц.
Управляя температурой или током накачки одного из лазеров при фиксированных
параметрах второго, можно реализовать ступенчатую перестройку частоты в
диапазоне 3{,}3--6{,}3\,ГГц (полоса $B = 3$\,ГГц), обеспечивающую
теоретическое пространственное разрешение в свободном пространстве
\begin{equation}
    \Delta r = \frac{c}{2 B} \approx
    \frac{3 \cdot 10^8}{2 \cdot 3 \cdot 10^9} = 0{,}05\;\text{м} = 5\;\text{см},
\end{equation}
что согласуется с выражением~\eqref{eq:range-resolution}.

\paragraph{Связь с измерительным циклом SFCW-радара.}
На каждом частотном шаге $f_k$ синхронный детектор выделяет компоненты $I_k$
и $Q_k$, которые после нормировки на амплитуду передаваемого сигнала образуют
комплексный S-параметр $S_{21}(f_k) \propto I_k + iQ_k$, как описано в
предыдущей главе. Набор значений $\{S_{21}(f_k)\}_{k=0}^{N-1}$ составляет
один полный свип и является исходными данными для всех дальнейших методов
обработки.

\paragraph{Основные параметры установки.}
В таблице~\ref{tab:system-params} сведены ключевые характеристики
экспериментальной установки.

\begin{table}[H]
\centering
\caption{Основные параметры экспериментальной установки}
\label{tab:system-params}
\begin{tabular}{l l}
\toprule
Параметр & Значение \\
\midrule
Диапазон рабочих частот & 3{,}3 -- 6{,}3\,ГГц \\
Полоса частот $B$ & 3\,ГГц \\
Длина волны лазеров & $\sim$1550\,нм (C-диапазон) \\
Точность стабилизации температуры & 0{,}05\,$^\circ$C \\
Разрешение по току & 0{,}002\,мА \\
Разрядность АЦП & 16\,бит \\
Частота внешнего тактового сигнала & 2\,МГц \\
Частота стробирования & 20\,кГц \\
Интерфейсы связи & UART 115200, USB 2.0 \\
Теоретическое разрешение по дальности & $\sim$5\,см (в воздухе) \\
\bottomrule
\end{tabular}
\end{table}
